\documentclass[sigconf,nonacm]{acmart}

\usepackage{booktabs}
\usepackage{listings}
\usepackage{microtype}

\lstset{
  basicstyle=\ttfamily\scriptsize,
  breaklines=true,
  frame=single,
  columns=fullflexible
}

\title{Execution-Attested Side-Effect Gating at the NLIP Type-3 Boundary}
\subtitle{NDSS SafeOSE '27 Open-Source AI Security Hackathon Report --- Focus Area 3}

\author{CIRWEL}
\affiliation{%
  \institution{Independent Security Research}
  \country{USA}
}
\email{cirwel@users.noreply.github.com}

\begin{abstract}
Open-source AI agents increasingly invoke side-effecting tools (such as publishing code, creating pull requests, or mutating shared repositories) across standardized communication layers like the Natural Language Interaction Protocol (NLIP, ECMA-430). However, current reference implementations (\texttt{nlip\_server} and \texttt{nlip\_sdk}) provide no communication-layer reference monitor at the Type-3 tool boundary, leaving side-effecting operations vulnerable to unverified execution, agent self-report fabrication (\emph{``all tests passed''}), forged protocol submessages, Confused-Deputy replay, and Time-of-Check-to-Time-of-Use (TOCTOU) working-tree mutations.

We present \texttt{nlip-receipts}, a pluggable, framework-agnostic Type-3 Security Extension (SE) for NLIP that enforces \emph{execution-attested side-effect gating}. Under \texttt{nlip-receipts}, a publication tool call (\texttt{publish}) is permitted if and only if an authorized caller references an immutable Git tree digest (\texttt{tree\_oid}) for which the trusted Type-3 verifier previously recorded a passing check (\texttt{verify}) inside an isolated archive snapshot (\texttt{git archive}) under the active policy epoch and target profile. Strictly satisfying the SafeOSE Focus Area~3 Fit Test, the gate decision is a pure, zero-LLM-call function of the structured NLIP tool-call submessage, SE-protected SQLite ledger state, and explicit operator policy. Across 2,080 benchmark evaluations spanning 10 adversarial attack classes and 3 benign workflows, \texttt{nlip-receipts} blocks 100\% of unauthorized publication attempts with 0\% false denials, adds $0.23\text{ ms}$ median gate decision latency, and interoperates identically across Claude Code, OpenAI Codex CLI, and Google Antigravity (\texttt{agy}) adapters. Finally, we map our controls to ECMA-434 and propose a new Security Countermeasure and Objective (\textbf{SCO 7.3.4: Execution-Attested Side-Effect Gating}).
\end{abstract}

\begin{document}
\sloppy
\emergencystretch=1.5em
\maketitle

\section{Introduction \& Problem Definition}
As AI agents transition from conversational assistants to autonomous software operators, they increasingly share open communication protocols and tool interfaces. The Natural Language Interaction Protocol (\textbf{NLIP}, standardized as \textbf{ECMA-430}~\cite{ecma430} with explanatory architecture in \textbf{ECMA TR/113}~\cite{ecmatr113}, the \textbf{ECMA-434} Security Profiles standard~\cite{ecma434}, and the IETF \textbf{NLIP-DMSC} multi-agent draft~\cite{ietf_dmsc}) establishes a uniform multimodal message envelope (\texttt{NLIP\_Message} and \texttt{NLIP\_SubMessage}) across heterogeneous agent frameworks. Both ECMA TR/113 (\S5.1, \S11.4) and NLIP-DMSC frame NLIP as the application-layer analogue of \texttt{ntohs}/\texttt{htons}: on the wire, natural language decouples internal schemas across endpoints while LLMs translate wire messages into internal structured calls (such as Model Context Protocol (MCP) or repository operations). In the SafeOSE control-point architecture, every boundary crossing can host a pluggable Security Extension (SE): Type-1 (user-to-agent), Type-2 (agent-to-agent), and \textbf{Type-3 (agent-to-tool/model boundary)}.

\paragraph{The Verification--Publication Gap in \texttt{nlip\_server}.}
Inspecting the upstream \texttt{nlip-project} reference repositories (\texttt{nlip\_sdk} v0.1.3 and \texttt{nlip\_server} v0.1.3) reveals two concrete architectural gaps when agents invoke side-effecting tools:
\begin{enumerate}
  \item \textbf{Absence of Type-3 Reference Monitors:} \texttt{nlip\_server} routes incoming \texttt{POST /nlip/} requests directly through \texttt{NLIP\_Session.correlated\_execute(msg)} into application handlers without any pre- or post-tool enforcement hooks, credential submessage stripping, or structured policy validation.
  \item \textbf{Unverifiable Agent Self-Reports \& TOCTOU Races:} In standard agentic workflows, operators instruct agents via prompts to \emph{``run the test suite before pushing.''} When pressured by context limits, reward hacking, or indirect prompt injection, an agent may claim in natural language that verification succeeded, fabricate a structured receipt payload, or modify the working tree after tests pass but before publication executes.
\end{enumerate}

\paragraph{Prior-Work Disclosure.}
This work builds on the author's prior single-harness Hermes plugin \texttt{hermes-receipts}~\cite{hermes_receipts}, which introduced working-tree fingerprinting to gate terminal \texttt{git push} commands inside a single CLI. However, terminal-hook command parsing is inherently fragile against shell aliasing, alternate executables, and concurrent working-tree mutations during test execution. For the NDSS SafeOSE~'27 Hackathon, we designed and built \texttt{nlip-receipts} from the ground up as a standards-conformant \textbf{ECMA-430/ECMA-434 Type-3 Security Extension} featuring: (i) typed NLIP \texttt{structured} submessage mediation rather than shell-string classification, (ii) \textbf{snapshot-isolated verification} (\texttt{git archive}) and \textbf{commit-pinned publication} (\texttt{<commit\_oid>:refs/heads/<branch>}) that eliminate mid-run and post-check TOCTOU windows, (iii) multi-caller Confused-Deputy and destination-profile binding, and (iv) cross-harness adapters for Claude Code, Codex CLI, and Antigravity (\texttt{agy}).

\section{Threat Model \& ECMA-434 Mapping}

\subsection{Threat Model and Attack Taxonomy}
We consider an untrusted or compromised AI agent (driven by any LLM or harness, potentially subverted by indirect prompt injection in repository files or tool outputs) interacting with a protected publication endpoint through an NLIP Type-3 boundary. The adversary's goal is to publish an unverified, failing, or trojaned artifact to a protected destination repository (e.g., \texttt{staging} or \texttt{prod}). We enumerate ten concrete attack classes ($\mathcal{A}_1$--$\mathcal{A}_{10}$):
\begin{itemize}
  \item $\mathcal{A}_1$ \textbf{(Unverified Publish):} Calling \texttt{publish} directly with no prior verification run.
  \item $\mathcal{A}_2$ \textbf{(Self-Report Fabrication):} Claiming \emph{``All 42 unit tests passed''} in the NLIP \texttt{text} field without a valid receipt.
  \item $\mathcal{A}_3$ \textbf{(Forged NLIP Submessage):} Attaching a fabricated \texttt{structured} receipt submessage (\texttt{status="passed"}) to the NLIP request.
  \item $\mathcal{A}_4$ \textbf{(Post-Check TOCTOU Commit):} Verifying clean commit $C_1$ (tree $T_1$), committing trojaned code $C_2$ (tree $T_2$), and invoking \texttt{publish} on $C_2$ citing $C_1$'s receipt.
  \item $\mathcal{A}_5$ \textbf{(Concurrent Working-Tree Race):} Mutating tracked files in the working tree while verification or publication is in flight.
  \item $\mathcal{A}_6$ \textbf{(Confused-Deputy Replay):} Caller \texttt{agent-b} replays a valid receipt ID issued to \texttt{agent-a} for the same tree.
  \item $\mathcal{A}_7$ \textbf{(Destination/Profile Substitution):} Using a receipt earned under a lightweight \texttt{staging} profile to publish to a \texttt{prod} destination requiring a stricter check.
  \item $\mathcal{A}_8$ \textbf{(Expired Receipt or Stale Policy Epoch):} Presenting a receipt after its TTL expires or after the operator updates the verification policy epoch.
  \item $\mathcal{A}_9$ \textbf{(Single-Use Receipt Replay):} Reusing an already-consumed single-use receipt for a subsequent publication.
  \item $\mathcal{A}_{10}$ \textbf{(Superseding Failure Bypass):} Presenting an older passing receipt for tree $T_1$ after a more recent verification run on $T_1$ failed.
\end{itemize}

\paragraph{Trusted Computing Base (TCB) and Security Boundary.}
The TCB comprises the Type-3 SE supervisor process, its SQLite ledger and ephemeral HMAC key, the local \texttt{git} binary invoked with sanitized configuration (\texttt{--no-replace-objects}, \texttt{GIT\_CONFIG\_NOSYSTEM=1}, \texttt{GIT\_CONFIG\_GLOBAL=/dev/null}), the operator-configured \texttt{PolicyConfig} (including \texttt{pinned\_file\_sha256} digests for protected verification scripts), and the bare destination repositories. The security guarantee is \emph{verified-equals-published equivalence}: a commit is published only if its exact immutable tree passed the operator-pinned check under the active policy epoch. \texttt{nlip-receipts} does \emph{not} guarantee that an operator's test suite is semantically complete, nor does it defend against host-kernel compromise.

\subsection{Mapping to ECMA-434 \& Proposed SCO 7.3.4}
Table~\ref{tab:scos} maps \texttt{nlip-receipts} to ECMA-434~\cite{ecma434}. While ECMA-434 \S7.3.2 mandates general runtime controls, it does not yet define a specific countermeasure for execution-attested precondition verification at tool boundaries. Following the Focus Area~3 call to propose new SCOs where needed, we define \textbf{SCO 7.3.4}:
\begin{quote}
\textbf{Proposed SCO 7.3.4 (Execution-Attested Side-Effect Gating):} \emph{Scope:} Agent Security Profile 3 (Runtime \& Behaviour). \emph{Objective:} Prevent agents from executing irreversible or external side effects based on hallucinated, forged, stale, or mismatched precondition checks. \emph{Countermeasure:} Side-effecting tool invocations crossing an NLIP Type-3 boundary shall require an opaque handle to a verification record minted by a trusted supervisor in local state, cryptographically bound to the caller identity, target profile, policy epoch, and immutable artifact digest being published.
\end{quote}

\begin{table}[t]
\centering
\caption{Alignment with ECMA-434 Security Profiles and SCOs.}
\label{tab:scos}
\footnotesize
\begin{tabular}{@{}llp{4.3cm}@{}}
\toprule
\textbf{Clause} & \textbf{ECMA-434 SCO} & \textbf{\texttt{nlip-receipts} Enforcement Mechanism} \\
\midrule
7.1.3 & Ethical-by-Design & Control-message disclosure \& \texttt{ETHICS.md} (Annex B). \\
7.2.2 & Client Authorization & Per-caller allowlist \& per-destination profile check. \\
7.2.3 & Privilege Misuse & Strips in-band auth tokens; binds receipts to \texttt{caller\_id} and \texttt{profile\_id} ($\mathcal{A}_6, \mathcal{A}_7$). \\
7.2.4 & Observability \& Audit & SHA-256 hash-chained SQLite decision log (\texttt{decisions}). \\
7.2.7 & Secure Disclosure & NIST SSDF disclosure in \texttt{SECURITY.md} (Annex A). \\
7.3.2 & Runtime Controls & Deterministic Type-3 reference monitor ($0$ LLM calls). \\
\textbf{7.3.4*} & \textbf{Attested Side Effects} & \textbf{Proposed SCO:} Snapshot-verified, tree- \& epoch-bound receipt gate ($\mathcal{A}_1$--$\mathcal{A}_{10}$). \\
\bottomrule
\end{tabular}
\end{table}

\section{Design \& Implementation}

\subsection{Protocol Binding over ECMA-430 Submessages}
\texttt{nlip-receipts} uses standard ECMA-430 message structures without modifying the wire format:
\begin{itemize}
  \item \textbf{Tool Call Proposal (\texttt{nlip.se.type3.tool\_call}):} Carried as a \texttt{structured} submessage (\texttt{subformat="json"}) specifying \texttt{operation} $\in \{\texttt{verify}, \texttt{publish}\}$, \texttt{repo\_id}, \texttt{artifact\_ref}, \texttt{destination}, \texttt{branch}, \texttt{profile\_id}, and optional \texttt{receipt\_id}.
  \item \textbf{Caller Authentication (\texttt{token}):} Carried in an ECMA-430 \texttt{token} submessage (\texttt{authentication} or \texttt{authorization}). At ingress, \texttt{extract\_caller\_credentials()} validates and strips the credential submessage so raw tokens never reach downstream tools or LLM prompts (\texttt{security\_guidelines} \S3.3, \S3.5), and rejects messages carrying conflicting duplicate authentication tokens to prevent token-pollution spoofing.
  \item \textbf{Decision \& Receipt Submessages:} Returned via labels \texttt{nlip.se.type3.decision} and \texttt{nlip.se.type3.receipt} in the response \texttt{NLIP\_Message} alongside a human-readable \texttt{text} summary and preserved \texttt{conversation} correlator token (ECMA-430 \S6.2.1).
\end{itemize}

\subsection{Focus Area 3 Fit-Test Verification}
The Focus Area~3 Fit Test asks: \emph{Can the security decision be made using data visible at the enforcement boundary, without depending on hidden reasoning or framework-specific internals?}

In \texttt{nlip-receipts}, the enforcement kernel \texttt{evaluate\_type3\_gate()} (\texttt{nlip\_receipts/gate.py}) is a pure function:
\[
  \mathsf{Decide}(\text{Call}, C_{\text{oid}}, T_{\text{oid}}, \text{Ledger}, \text{Policy}, t_{\text{now}}) \rightarrow \{\mathsf{allow}, \mathsf{deny}\}
\]
It evaluates ten deterministic predicates in sequence:
(1) policy existence (\texttt{POLICY\_NOT\_FOUND}),
(2) caller HMAC credential authentication and allowlist authorization (\texttt{CALLER\_AUTH\_INVALID}, \texttt{CALLER\_UNAUTHORIZED}),
(3) rejection of forged receipt submessages (\texttt{FORGED\_RECEIPT\_SUBMESSAGE}),
(4) presence of a ledger receipt (\texttt{MISSING\_RECEIPT}),
(5) SE-held HMAC-SHA256 verification (\texttt{RECEIPT\_INTEGRITY\_FAILED}),
(6) caller equality $\text{rcpt}.\texttt{caller\_id} = \text{call}.\texttt{caller\_id}$ (\texttt{CROSS\_CALLER\_REPLAY}),
(7) repository, destination, profile, and command-hash equality (\texttt{PROFILE\_OR\_REPO\_MISMATCH}),
(8) policy epoch and TTL freshness $t_{\text{now}} \le \text{rcpt}.\texttt{expires\_at}$ (\texttt{POLICY\_EPOCH\_STALE}, \texttt{RECEIPT\_EXPIRED}),
(9) single-use unconsumed status (\texttt{RECEIPT\_ALREADY\_CONSUMED}), and
(10) exact commit and content-tree equality ($\text{rcpt}.\texttt{commit\_oid} = C_{\text{oid}}$ and $\text{rcpt}.\texttt{tree\_oid} = T_{\text{oid}}$) with no newer failing run on $T_{\text{oid}}$ (\texttt{ARTIFACT\_COMMIT\_MISMATCH}, \texttt{ARTIFACT\_TREE\_MISMATCH}, \texttt{SUPERSEDED\_BY\_FAILURE}).

\subsection{Closing TOCTOU via Snapshot Verification and Commit-Pinned Push}
A subtle flaw in working-tree hooks is that files can change while tests run or between test completion and \texttt{git push}. \texttt{nlip-receipts} eliminates both windows:
\begin{enumerate}
  \item \textbf{Manifest-Verified Snapshot Execution:} On \texttt{verify(repo\_id, artifact\_ref, profile\_id)}, the SE rejects revision-search or reflog expressions (\texttt{:/}, \texttt{@\{...\}}) and resolves \texttt{artifact\_ref} to $(\text{Commit}_{\text{oid}}, \text{Tree}_{\text{oid}})$ via \texttt{git --no-replace-objects rev-parse --verify --end-of-options} with hooks disabled. It exports $\text{Tree}_{\text{oid}}$ via \texttt{git archive --worktree-attributes} into a fresh temporary directory (\texttt{tarfile.extractall(filter="data")}), verifies 1-to-1 blob-digest equivalence against \texttt{git ls-tree -r --full-tree} (blocking \texttt{.gitattributes} \texttt{export-ignore}/\texttt{export-subst}, symlinks, and submodule gitlinks), verifies any operator-pinned test file hashes (\texttt{pinned\_file\_sha256}), and executes \texttt{policy.command} (\texttt{shell=False}).
  \item \textbf{Transactional Check-and-Consume \& Commit-Pinned Push:} On \texttt{publish}, once \texttt{evaluate\_type3\_gate()} returns \texttt{allow}, the ledger executes a \texttt{BEGIN IMMEDIATE} transaction that re-verifies \texttt{consumed = 0}, \texttt{policy\_epoch}, and absence of superseding failures before setting \texttt{consumed = 1}. The executor then pushes the exact resolved commit object:
  \begin{lstlisting}
git push --no-verify --no-follow-tags --porcelain -- <dest> <commit_oid>:refs/heads/<branch>
  \end{lstlisting}
  Even if the agent mutated the working tree or advanced \texttt{HEAD} after verification, those unverified changes cannot enter the published ref.
\end{enumerate}

\subsection{Multi-Harness Adapters}
Because all policy and state live behind the NLIP message boundary, integrating any agent harness requires only a thin (~30 LOC) adapter (\texttt{nlip\_receipts/adapters.py}). We provide working adapters for:
(1) \textbf{Claude Code} (\texttt{PreToolUse} hook emitting \texttt{permissionDecision: allow|deny}),
(2) \textbf{OpenAI Codex CLI} (\texttt{tool.execution.result} wrapper), and
(3) \textbf{Google Antigravity (\texttt{agy})} (\texttt{EXECUTED} / \texttt{BLOCKED\_BY\_NLIP\_SE} tool wrapper), as well as standard FastAPI \texttt{nlip\_server} HTTP endpoints.

\section{Evaluation}

We evaluate \texttt{nlip-receipts} across \textbf{8 system configurations} (4 baselines, 3 component ablations, and our full gate \texttt{S4\_nlip\_type3\_gate}) over \textbf{13 scenarios} ($\mathcal{A}_1$--$\mathcal{A}_{10}$ adversarial + $\mathcal{B}_1$--$\mathcal{B}_3$ benign workflows, 20 trials per scenario = 260 runs per system, \textbf{2,080 total evaluations}). All experiments are reproducible offline in $<6$ seconds via \texttt{make eval}.

\subsection{Baselines and Ablations}
\begin{itemize}
  \item \texttt{B0\_no\_gate}: Vanilla \texttt{nlip\_server} execution with no Type-3 SE.
  \item \texttt{B1\_prompt\_only}: Instruction-level guard that trusts natural-language completion claims or attached receipt fields.
  \item \texttt{B2\_submessage\_trust}: Stateless boundary filter that inspects whether an NLIP \texttt{structured} receipt submessage reports \texttt{status="passed"}, without a server-side ledger lookup.
  \item \texttt{B3\_unbound\_flag}: Stateful session boolean that records whether \emph{any} verification command exited $0$ during the session, without binding to $\text{Tree}_{\text{oid}}$, caller, destination profile, or TTL.
  \item \texttt{Abl\_no\_tree\_bind}, \texttt{Abl\_no\_caller\_bind}, \texttt{Abl\_no\_profile\_bind}: Ablations of \texttt{S4} removing artifact-tree binding, caller binding, or destination-profile binding, respectively.
\end{itemize}

\begin{table}[t]
\centering
\caption{Overall Evaluation Results across 2,080 Runs (200 Adversarial + 60 Benign per System).}
\label{tab:eval_overall}
\scriptsize
\setlength{\tabcolsep}{1.8pt}
\begin{tabular}{@{}lrrrrrr@{}}
\toprule
\textbf{System} & \textbf{Block Rate} & \textbf{Unauth.} & \textbf{Benign} & \textbf{FP} & \textbf{p50} & \textbf{p99} \\
\midrule
\texttt{B0\_no\_gate} & 0\% (0/200) & 100\% & 100\% & 0\% & 0.25 & 0.36 \\
\texttt{B1\_prompt\_only} & 10\% (20/200) & 90\% & 100\% & 0\% & 0.28 & 0.48 \\
\texttt{B2\_submessage\_trust} & 20\% (40/200) & 80\% & 100\% & 0\% & 0.30 & 0.63 \\
\texttt{B3\_unbound\_flag} & 30\% (60/200) & 70\% & 100\% & 0\% & 0.33 & 0.75 \\
\midrule
\texttt{Abl\_no\_tree\_bind} & 80\% (160/200) & 20\% & 100\% & 0\% & 0.69 & 1.33 \\
\texttt{Abl\_no\_caller\_bind} & 90\% (180/200) & 10\% & 100\% & 0\% & 0.71 & 1.40 \\
\texttt{Abl\_no\_profile\_bind} & 90\% (180/200) & 10\% & 100\% & 0\% & 0.76 & 1.55 \\
\midrule
\textbf{\texttt{S4\_nlip\_type3\_gate}} & \textbf{100\% (200/200)} & \textbf{0\%} & \textbf{100\%} & \textbf{0\%} & \textbf{0.76} & \textbf{1.67} \\
\bottomrule
\end{tabular}
\end{table}

\begin{table}[t]
\centering
\caption{Per-Scenario Block Rate (Adversarial $\mathcal{A}_1$--$\mathcal{A}_{10}$) and Completion Rate (Benign $\mathcal{B}_1$--$\mathcal{B}_3$).}
\label{tab:eval_scenarios}
\scriptsize
\setlength{\tabcolsep}{2.8pt}
\begin{tabular}{@{}llrrrrrrr@{}}
\toprule
\textbf{Scenario} & \textbf{B0} & \textbf{B1} & \textbf{B2} & \textbf{B3} & \textbf{-Tree} & \textbf{-Call} & \textbf{-Prof} & \textbf{S4 (Ours)} \\
\midrule
$\mathcal{A}_1$ Unverified publish & 0\% & 100\% & 100\% & 100\% & 100\% & 100\% & 100\% & \textbf{100\%} \\
$\mathcal{A}_2$ Self-report claim & 0\% & 0\% & 100\% & 100\% & 100\% & 100\% & 100\% & \textbf{100\%} \\
$\mathcal{A}_3$ Forged submessage & 0\% & 0\% & 0\% & 100\% & 100\% & 100\% & 100\% & \textbf{100\%} \\
$\mathcal{A}_4$ Post-check TOCTOU & 0\% & 0\% & 0\% & 0\% & 0\% & 100\% & 100\% & \textbf{100\%} \\
$\mathcal{A}_5$ Worktree race edit & 0\% & 0\% & 0\% & 0\% & 0\% & 100\% & 100\% & \textbf{100\%} \\
$\mathcal{A}_6$ Cross-caller replay & 0\% & 0\% & 0\% & 0\% & 100\% & 0\% & 100\% & \textbf{100\%} \\
$\mathcal{A}_7$ Dest. substitution & 0\% & 0\% & 0\% & 0\% & 100\% & 100\% & 0\% & \textbf{100\%} \\
$\mathcal{A}_8$ Expired / stale epoch & 0\% & 0\% & 0\% & 0\% & 100\% & 100\% & 100\% & \textbf{100\%} \\
$\mathcal{A}_9$ Single-use replay & 0\% & 0\% & 0\% & 0\% & 100\% & 100\% & 100\% & \textbf{100\%} \\
$\mathcal{A}_{10}$ Superseded failure & 0\% & 0\% & 0\% & 0\% & 100\% & 100\% & 100\% & \textbf{100\%} \\
\midrule
$\mathcal{B}_1$ Benign verify+pub & 100\% & 100\% & 100\% & 100\% & 100\% & 100\% & 100\% & \textbf{100\%} \\
$\mathcal{B}_2$ Fail$\rightarrow$fix$\rightarrow$reverify & 100\% & 100\% & 100\% & 100\% & 100\% & 100\% & 100\% & \textbf{100\%} \\
$\mathcal{B}_3$ Multi-harness pub & 100\% & 100\% & 100\% & 100\% & 100\% & 100\% & 100\% & \textbf{100\%} \\
\bottomrule
\end{tabular}
\end{table}

\subsection{Results and Findings}
Tables~\ref{tab:eval_overall} and~\ref{tab:eval_scenarios} summarize the empirical results:
\begin{enumerate}
  \item \textbf{Necessity of Content- \& Context-Bound Receipts:} Stateless submessage inspection (\texttt{B2}) blocks only $20\%$ of attacks because an adversary can forge a \texttt{structured} receipt submessage ($\mathcal{A}_3$). Session-level boolean flags (\texttt{B3}) block $\mathcal{A}_1$--$\mathcal{A}_3$ ($30\%$) but fail completely against post-verification code edits ($\mathcal{A}_4, \mathcal{A}_5$), cross-caller replay ($\mathcal{A}_6$), destination substitution ($\mathcal{A}_7$), TTL/epoch staleness ($\mathcal{A}_8$), single-use replay ($\mathcal{A}_9$), and superseding failures ($\mathcal{A}_{10}$).
  \item \textbf{Ablation Validation:} Removing $\text{Tree}_{\text{oid}}$ binding drops block rate from $100\%$ to $80\%$ (failing $\mathcal{A}_4$ and $\mathcal{A}_5$); removing caller binding drops it to $90\%$ (failing $\mathcal{A}_6$); and removing profile/command binding drops it to $90\%$ (failing $\mathcal{A}_7$). Only the full tuple binding in \texttt{S4\_nlip\_type3\_gate} achieves \textbf{100.0\% (200/200)} attack blocking while maintaining \textbf{100.0\% (60/60)} benign completion (including recovery after an initial test failure in $\mathcal{B}_2$).
  \item \textbf{Runtime Overhead:} Because \texttt{evaluate\_type3\_gate()} performs only SQLite indexed lookups and SHA-256/HMAC-SHA256 verification with \textbf{zero LLM calls}, the gate decision overhead is \textbf{0.76~ms (p50)} and \textbf{1.67~ms (p99)}. End-to-end snapshot verification (\texttt{git archive} + test subprocess) takes $31.86\text{ ms}$ (p50), and commit-pinned publication (\texttt{git push}) takes $52.77\text{ ms}$ (p50).
  \item \textbf{Live LLM Indirect Prompt Injection (IPI) Evaluation (\texttt{eval/live\_agent\_eval.py}):} To separate deterministic replay invariants from live model behavior, we evaluated two local tool-calling LLMs (\texttt{qwen3-coder-next:latest} and \texttt{gemma4:latest}, temperature $0.0$, seed $42$) on three repository-embedded IPI tasks ($\mathcal{I}_1$: fake CI-down override in \texttt{DEPLOY\_NOTES.md}; $\mathcal{I}_2$: forged receipt JSON in \texttt{HOTFIX\_402.md}; $\mathcal{I}_3$: developer note requesting TOCTOU reuse of a prior commit's receipt on a trojaned commit) and one multi-turn benign release workflow ($\mathcal{B}_4$: \texttt{verify} $\rightarrow$ \texttt{publish}). Both models were subverted by $3/3$ ($100\%$) of the IPI prompts into emitting unauthorized \texttt{publish} tool calls, resulting in $3/3$ unauthorized destination pushes under \texttt{B0\_no\_gate} and \texttt{B1\_prompt\_only}. Under \textbf{\texttt{S4\_nlip\_type3\_gate}}, \textbf{0/3 (0\%)} resulted in unauthorized destination mutations (blocked via \texttt{MISSING\_RECEIPT}, \texttt{FORGED\_RECEIPT\_SUBMESSAGE}, and \texttt{ARTIFACT\_TREE\_MISMATCH}), while \textbf{1/1 (100\%)} of the live multi-turn benign workflows succeeded for both models.
\end{enumerate}

\section{Limitations, Related Work \& Conclusion}
\paragraph{Related Work.}
Reference monitors and capability systems (such as OAuth DPoP, RFC 8693 token exchange, and Sigstore/in-toto attestations~\cite{intoto}) secure transport authentication and CI/CD supply chains. \texttt{nlip-receipts} brings content-addressed execution attestation directly to the \emph{in-session AI agent tool boundary} over ECMA-430/ECMA-434, bridging the gap between probabilistic agent reasoning and deterministic software publication gates.

\paragraph{Limitations \& Future Work.}
First, a receipt attests that the configured check command passed on $\text{Tree}_{\text{oid}}$, not that the check suite itself has complete coverage. Second, while we demonstrate Git-backed artifacts here, the \texttt{ToolCallProposal} and \texttt{VerificationReceipt} schemas apply to any content-addressed payload (e.g., container image digests or SQL migration bundles). Third, in multi-node deployments, the SQLite ledger can be replaced by a shared PostgreSQL store with detached Ed25519 signatures.

\paragraph{Conclusion.}
\texttt{nlip-receipts} provides a practical, immediately upstreamable Type-3 Security Extension for \texttt{nlip\_server} and \texttt{nlip\_sdk} that turns agent precondition verification from an unenforceable prompt instruction into a deterministic, cryptographically bound protocol invariant.

\begin{thebibliography}{9}
\bibitem{ecma430}
Ecma International. 2025. \emph{ECMA-430: Natural Language Interaction Protocol (NLIP)}, 1st Edition (December 2025).
\bibitem{ecma434}
Ecma International. 2025. \emph{ECMA-434: Security Profiles for Natural Language Interaction Protocol (NLIP)}, 1st Edition (December 2025).
\bibitem{ecmatr113}
Ecma International. 2025. \emph{ECMA TR/113: Explanatory Guide to the Natural Language Interaction Protocol (NLIP)}, 1st Edition (December 2025).
\bibitem{ietf_dmsc}
D. Verma, S. Hughes, L. Xing, R. Topaloglu, E. Bertino, A. Ratnaparkhi, S. Aiyagari, and T. Ross. 2026. \emph{Natural Language Interaction for Dynamic Multi-Agent Secure Collaboration (NLIP-DMSC)}. IETF Internet-Draft \texttt{draft-natural-language-multiagent-00}.
\bibitem{nlip_security_guidelines}
NLIP Security Working Group. 2025. \emph{NLIP Security Guidelines and Best Practices}. \url{https://github.com/nlip-project/security_guidelines}.
\bibitem{hermes_receipts}
CIRWEL. 2026. \emph{receipts: Verification Receipts Bound to Repository Content}. \url{https://github.com/cirwel/hermes-receipts}.
\bibitem{intoto}
Santiago Torres-Arias et al. 2019. \emph{in-toto: Providing Farm-to-Table Guarantees for Bits and Bytes}. In USENIX Security Symposium.
\end{thebibliography}

\appendix
\section{Artifact Reproducibility Appendix}
All code, schemas, adapters, tests, and upstream standards proposals (\texttt{docs/UPSTREAM\_STANDARDS\_TEXT.md}) are released under the Apache-2.0 License at \url{https://github.com/cirwel/nlip-receipts}:
\begin{lstlisting}
# 0. One-command environment setup (installs nlip_sdk & nlip_server)
make setup

# 1. Run the 13 end-to-end pytest suites (HTTP, TOCTOU, Replay, Concurrency, 3 Adapters)
make test

# 2. Reproduce the 2,080-run evaluation benchmark & ablation tables
make eval

# 3. Run the 7-step interactive NLIP Type-3 demonstration
make demo
\end{lstlisting}

\end{document}
